\documentclass[12pt]{article}
\usepackage[utf8]{inputenc}
\usepackage[T1]{fontenc}
\usepackage{hyperref}
\title{Darkinyung Language Resources for K-12 Curriculum Development}
\author{Compiled by Claude Sonnet 5.5 (claude-sonnet-5-5) for Andr\'as Kornai}
\date{}
\begin{document}
\maketitle

\section{Introduction}
Darkinyung (also spelled Darkinjung or Darkinung; ISO 639-3 xda; Glottocode hawk1239) is a Pama-Nyungan language of the Central Coast and Hawkesbury region of New South Wales. It is no longer spoken as a first language. A revival effort exists, and the Darkinyung Language Group was formed in 2003. Its sources are historical, mainly Robert Mathews' No.\ 7 notebook, held at the National Library of Australia. No audio recordings survive.

\section{Primary Educational Resources}
I found no school curriculum or classroom materials specific to Darkinyung. The grammar and dictionary below are the nearest thing to a teaching resource.

\section{Linguistic and Reference Materials}
\emph{Darkinyung Grammar and Dictionary: Revitalising a Language from Historical Sources}, by Caroline Jones, is published by Muurrbay Aboriginal Language and Culture Co-operative: \url{https://muurrbay.org.au/publications/dictionaries-grammar/darkinyung-grammar-and-dictionary-revitalising-a-language-from-historical-sources/}. Because the historical wordlists are sparse, the team supplemented them with words from lexically similar neighbouring languages. A thesis abstract on the language is at \url{https://ses.library.usyd.edu.au/bitstream/handle/2123/7745/GE%20Ford%202013%2c%20Darkinung%20Abstract.pdf}. Overviews are at \url{https://en.wikipedia.org/wiki/Darkinjung_language} and \url{https://historyofaboriginalsydney.edu.au/north-west/darkinung-country-darkinung-people-and-darkinung-culture}. The 1829 list of Darug and Darkinyung place names along the Hawkesbury River by Reverend John McGarvie is held at the State Library of NSW. It has been used in the Dyarubbin Project, but I saw no URL for it.

\section{Religious and Community Resources}
The Darkinyung Language Group is the community body, but I did not see a website for it. The Muurrbay Co-operative is the publisher of the grammar and dictionary.

\section{Digital and Audio Resources}
The PARADISEC blog post on the launch of the book is at \url{https://www.paradisec.org.au/blog/2008/12/brand-new-day-for-the-darkinyung-language}. Dictionary of Sydney entries that came up in the search are \url{https://dictionaryofsydney.org/node/67399}, \url{https://dictionaryofsydney.org/node/65937} and \url{https://dictionaryofsydney.org/node/10577}. I did not read them closely and cannot say how much language content they hold. No audio resources were found.

\section{Curriculum Development Context}
The Jones grammar and dictionary give a basis for introducing the language in schools, with Darkinyung vocabulary, place names and cultural content. Because the material is reconstructed from 19th-century records and supplemented from neighbouring languages, any curriculum should be developed with the Darkinyung community and language group. I found no school currently teaching it.
\end{document}